%=====================================================================
%   APPENDIX D: TOPIC DESCRIPTION
%=====================================================================
%   The topic description as it was issued during project selection, before
%   the problem was refined.
%
%   Required for a final year project report. Not usually included in a
%   postgraduate thesis: comment out %=====================================================================
%   APPENDIX D: TOPIC DESCRIPTION
%=====================================================================
%   The topic description as it was issued during project selection, before
%   the problem was refined.
%
%   Required for a final year project report. Not usually included in a
%   postgraduate thesis: comment out %=====================================================================
%   APPENDIX D: TOPIC DESCRIPTION
%=====================================================================
%   The topic description as it was issued during project selection, before
%   the problem was refined.
%
%   Required for a final year project report. Not usually included in a
%   postgraduate thesis: comment out %=====================================================================
%   APPENDIX D: TOPIC DESCRIPTION
%=====================================================================
%   The topic description as it was issued during project selection, before
%   the problem was refined.
%
%   Required for a final year project report. Not usually included in a
%   postgraduate thesis: comment out \include{C_topicdescription} in
%   0_main.tex if you are writing one.
%
%   HOW TO USE THIS
%     1. Put the PDF of the topic list in pdfs/.
%     2. Uncomment the \includepdf below and set the page range to YOUR topic.
%
%   Nothing is included by default, because the PDF is specific to a course
%   and a year and does not belong in a template. Leaving a hardcoded
%   \includepdf here would break the build for everyone who cloned the repo
%   without that file.
%=====================================================================

\chapter*{Appendix D: Topic description}
\addcontentsline{toc}{chapter}{Appendix D: Topic description}
%---------------------------------------------------------------------
The original topic description used during the project selection process is
included below. This topic has since been further refined, and the exact
problem that was investigated is described in
Section~\ref{sec:ch1.problemStatement}.

%---------------------------------------------------------------------
%   Subject of investigation, quoted from the EEE4022S topic list
%   (RA Verrinder, 8 July 2026). Its references are listed below by hand,
%   not through BibTeX, so they stay separate from the main bibliography.
%---------------------------------------------------------------------
\section*{Subject of investigation}
\begin{quote}
Building on Zhao et al.\ (2022) and the surface characterisation work of Spirakis (2021); Spirakis et al.\ (2023);
Jones (2024), this project develops and validates a stereo-camera three-dimensional reconstruction pipeline on the
UCT BlueROV2 platform, as a first step towards reconstructing the underside of Antarctic sea ice. The work is
conducted in a controlled pool environment, where a submerged proxy target of known geometry stands in for the ice
underside. External fiducial markers provide the camera pose reference, so that the reconstruction is decoupled from
onboard navigation, and the recovered surface model is evaluated against the ground-truth geometry of the proxy. The
focus is on the geometric reconstruction itself: stereo calibration, handling of the underwater optical path,
multi-view fusion, and quantified reconstruction accuracy.
\end{quote}

\section*{References}
\begin{list}{}{\setlength{\leftmargin}{1.5em}\setlength{\itemindent}{-1.5em}\setlength{\itemsep}{0.3em}}
\item Blue Robotics. BlueROV2. [ONLINE], 2025. Accessed: 2025-07-11.
  url: \url{https://bluerobotics.com/store/rov/bluerov2/}
\item D. Jones. SHARC project: Ice floe size and shape characterisation from LiDAR scans. BSc(Eng) final year
  project, Department of Electrical Engineering), University of Cape Town, 2024.
  url: \url{https://tinyurl.com/djones2023}
\item J. Landy, A. Komarov, and D. Barber. Numerical and experimental evaluation of terrestrial LiDAR for
  parameterizing centimeter-scale sea ice surface roughness. \textit{IEEE Transactions on Geoscience and Remote
  Sensing}, 53:4887--4898, 2015. doi: \href{https://doi.org/10.1109/TGRS.2015.2412034}{10.1109/TGRS.2015.2412034}
\item J. Landy, A. Komarov, and D. Barber. Secondary-scale surface roughness parameterization using terrestrial
  LiDAR. In \textit{IGARSS 2015-2015 IEEE International Geoscience and Remote Sensing Symposium}, pages
  2068--2017. IEEE, 2023. doi: \href{https://doi.org/10.1109/igarss.2015.7326208}{10.1109/igarss.2015.7326208}
\item A. Spirakis. Investigation into surface roughness characterisation using LiDAR and cameras. BSc(Eng) final
  year project, Department of Electrical Engineering), University of Cape Town, 2021.
  url: \url{https://tinyurl.com/spirakis2021}
\item A. Spirakis, J. Hepworth, and R. Verrinder. 3D reconstruction of pancake sea ice using LiDAR and cameras. In
  \textit{IGARSS 2023-2023 IEEE International Geoscience and Remote Sensing Symposium}, pages 60--63. IEEE, 2023.
  doi: \href{https://doi.org/10.1109/igarss52108.2023.10282250}{10.1109/igarss52108.2023.10282250}
\item M. Vichi. An indicator of sea ice variability for the Antarctic marginal ice zone. \textit{The Cryosphere},
  16:4087--4106, 2022. doi: \href{https://doi.org/10.5194/tc-16-4087-2022}{10.5194/tc-16-4087-2022}
\item L. Zhao, M. Zhou, and B. Loose. Towards under-ice sensing using a portable ROV. In \textit{OCEANS
  2022-Hampton Roads}, pages 1--8. IEEE, 2022.
  doi: \href{https://doi.org/10.1109/OCEANS47191.2022.9977140}{10.1109/OCEANS47191.2022.9977140}
\end{list}
 in
%   0_main.tex if you are writing one.
%
%   HOW TO USE THIS
%     1. Put the PDF of the topic list in pdfs/.
%     2. Uncomment the \includepdf below and set the page range to YOUR topic.
%
%   Nothing is included by default, because the PDF is specific to a course
%   and a year and does not belong in a template. Leaving a hardcoded
%   \includepdf here would break the build for everyone who cloned the repo
%   without that file.
%=====================================================================

\chapter*{Appendix D: Topic description}
\addcontentsline{toc}{chapter}{Appendix D: Topic description}
%---------------------------------------------------------------------
The original topic description used during the project selection process is
included below. This topic has since been further refined, and the exact
problem that was investigated is described in
Section~\ref{sec:ch1.problemStatement}.

%---------------------------------------------------------------------
%   Subject of investigation, quoted from the EEE4022S topic list
%   (RA Verrinder, 8 July 2026). Its references are listed below by hand,
%   not through BibTeX, so they stay separate from the main bibliography.
%---------------------------------------------------------------------
\section*{Subject of investigation}
\begin{quote}
Building on Zhao et al.\ (2022) and the surface characterisation work of Spirakis (2021); Spirakis et al.\ (2023);
Jones (2024), this project develops and validates a stereo-camera three-dimensional reconstruction pipeline on the
UCT BlueROV2 platform, as a first step towards reconstructing the underside of Antarctic sea ice. The work is
conducted in a controlled pool environment, where a submerged proxy target of known geometry stands in for the ice
underside. External fiducial markers provide the camera pose reference, so that the reconstruction is decoupled from
onboard navigation, and the recovered surface model is evaluated against the ground-truth geometry of the proxy. The
focus is on the geometric reconstruction itself: stereo calibration, handling of the underwater optical path,
multi-view fusion, and quantified reconstruction accuracy.
\end{quote}

\section*{References}
\begin{list}{}{\setlength{\leftmargin}{1.5em}\setlength{\itemindent}{-1.5em}\setlength{\itemsep}{0.3em}}
\item Blue Robotics. BlueROV2. [ONLINE], 2025. Accessed: 2025-07-11.
  url: \url{https://bluerobotics.com/store/rov/bluerov2/}
\item D. Jones. SHARC project: Ice floe size and shape characterisation from LiDAR scans. BSc(Eng) final year
  project, Department of Electrical Engineering), University of Cape Town, 2024.
  url: \url{https://tinyurl.com/djones2023}
\item J. Landy, A. Komarov, and D. Barber. Numerical and experimental evaluation of terrestrial LiDAR for
  parameterizing centimeter-scale sea ice surface roughness. \textit{IEEE Transactions on Geoscience and Remote
  Sensing}, 53:4887--4898, 2015. doi: \href{https://doi.org/10.1109/TGRS.2015.2412034}{10.1109/TGRS.2015.2412034}
\item J. Landy, A. Komarov, and D. Barber. Secondary-scale surface roughness parameterization using terrestrial
  LiDAR. In \textit{IGARSS 2015-2015 IEEE International Geoscience and Remote Sensing Symposium}, pages
  2068--2017. IEEE, 2023. doi: \href{https://doi.org/10.1109/igarss.2015.7326208}{10.1109/igarss.2015.7326208}
\item A. Spirakis. Investigation into surface roughness characterisation using LiDAR and cameras. BSc(Eng) final
  year project, Department of Electrical Engineering), University of Cape Town, 2021.
  url: \url{https://tinyurl.com/spirakis2021}
\item A. Spirakis, J. Hepworth, and R. Verrinder. 3D reconstruction of pancake sea ice using LiDAR and cameras. In
  \textit{IGARSS 2023-2023 IEEE International Geoscience and Remote Sensing Symposium}, pages 60--63. IEEE, 2023.
  doi: \href{https://doi.org/10.1109/igarss52108.2023.10282250}{10.1109/igarss52108.2023.10282250}
\item M. Vichi. An indicator of sea ice variability for the Antarctic marginal ice zone. \textit{The Cryosphere},
  16:4087--4106, 2022. doi: \href{https://doi.org/10.5194/tc-16-4087-2022}{10.5194/tc-16-4087-2022}
\item L. Zhao, M. Zhou, and B. Loose. Towards under-ice sensing using a portable ROV. In \textit{OCEANS
  2022-Hampton Roads}, pages 1--8. IEEE, 2022.
  doi: \href{https://doi.org/10.1109/OCEANS47191.2022.9977140}{10.1109/OCEANS47191.2022.9977140}
\end{list}
 in
%   0_main.tex if you are writing one.
%
%   HOW TO USE THIS
%     1. Put the PDF of the topic list in pdfs/.
%     2. Uncomment the \includepdf below and set the page range to YOUR topic.
%
%   Nothing is included by default, because the PDF is specific to a course
%   and a year and does not belong in a template. Leaving a hardcoded
%   \includepdf here would break the build for everyone who cloned the repo
%   without that file.
%=====================================================================

\chapter*{Appendix D: Topic description}
\addcontentsline{toc}{chapter}{Appendix D: Topic description}
%---------------------------------------------------------------------
The original topic description used during the project selection process is
included below. This topic has since been further refined, and the exact
problem that was investigated is described in
Section~\ref{sec:ch1.problemStatement}.

%---------------------------------------------------------------------
%   Subject of investigation, quoted from the EEE4022S topic list
%   (RA Verrinder, 8 July 2026). Its references are listed below by hand,
%   not through BibTeX, so they stay separate from the main bibliography.
%---------------------------------------------------------------------
\section*{Subject of investigation}
\begin{quote}
Building on Zhao et al.\ (2022) and the surface characterisation work of Spirakis (2021); Spirakis et al.\ (2023);
Jones (2024), this project develops and validates a stereo-camera three-dimensional reconstruction pipeline on the
UCT BlueROV2 platform, as a first step towards reconstructing the underside of Antarctic sea ice. The work is
conducted in a controlled pool environment, where a submerged proxy target of known geometry stands in for the ice
underside. External fiducial markers provide the camera pose reference, so that the reconstruction is decoupled from
onboard navigation, and the recovered surface model is evaluated against the ground-truth geometry of the proxy. The
focus is on the geometric reconstruction itself: stereo calibration, handling of the underwater optical path,
multi-view fusion, and quantified reconstruction accuracy.
\end{quote}

\section*{References}
\begin{list}{}{\setlength{\leftmargin}{1.5em}\setlength{\itemindent}{-1.5em}\setlength{\itemsep}{0.3em}}
\item Blue Robotics. BlueROV2. [ONLINE], 2025. Accessed: 2025-07-11.
  url: \url{https://bluerobotics.com/store/rov/bluerov2/}
\item D. Jones. SHARC project: Ice floe size and shape characterisation from LiDAR scans. BSc(Eng) final year
  project, Department of Electrical Engineering), University of Cape Town, 2024.
  url: \url{https://tinyurl.com/djones2023}
\item J. Landy, A. Komarov, and D. Barber. Numerical and experimental evaluation of terrestrial LiDAR for
  parameterizing centimeter-scale sea ice surface roughness. \textit{IEEE Transactions on Geoscience and Remote
  Sensing}, 53:4887--4898, 2015. doi: \href{https://doi.org/10.1109/TGRS.2015.2412034}{10.1109/TGRS.2015.2412034}
\item J. Landy, A. Komarov, and D. Barber. Secondary-scale surface roughness parameterization using terrestrial
  LiDAR. In \textit{IGARSS 2015-2015 IEEE International Geoscience and Remote Sensing Symposium}, pages
  2068--2017. IEEE, 2023. doi: \href{https://doi.org/10.1109/igarss.2015.7326208}{10.1109/igarss.2015.7326208}
\item A. Spirakis. Investigation into surface roughness characterisation using LiDAR and cameras. BSc(Eng) final
  year project, Department of Electrical Engineering), University of Cape Town, 2021.
  url: \url{https://tinyurl.com/spirakis2021}
\item A. Spirakis, J. Hepworth, and R. Verrinder. 3D reconstruction of pancake sea ice using LiDAR and cameras. In
  \textit{IGARSS 2023-2023 IEEE International Geoscience and Remote Sensing Symposium}, pages 60--63. IEEE, 2023.
  doi: \href{https://doi.org/10.1109/igarss52108.2023.10282250}{10.1109/igarss52108.2023.10282250}
\item M. Vichi. An indicator of sea ice variability for the Antarctic marginal ice zone. \textit{The Cryosphere},
  16:4087--4106, 2022. doi: \href{https://doi.org/10.5194/tc-16-4087-2022}{10.5194/tc-16-4087-2022}
\item L. Zhao, M. Zhou, and B. Loose. Towards under-ice sensing using a portable ROV. In \textit{OCEANS
  2022-Hampton Roads}, pages 1--8. IEEE, 2022.
  doi: \href{https://doi.org/10.1109/OCEANS47191.2022.9977140}{10.1109/OCEANS47191.2022.9977140}
\end{list}
 in
%   0_main.tex if you are writing one.
%
%   HOW TO USE THIS
%     1. Put the PDF of the topic list in pdfs/.
%     2. Uncomment the \includepdf below and set the page range to YOUR topic.
%
%   Nothing is included by default, because the PDF is specific to a course
%   and a year and does not belong in a template. Leaving a hardcoded
%   \includepdf here would break the build for everyone who cloned the repo
%   without that file.
%=====================================================================

\chapter*{Appendix D: Topic description}
\addcontentsline{toc}{chapter}{Appendix D: Topic description}
%---------------------------------------------------------------------
The original topic description used during the project selection process is
included below. This topic has since been further refined, and the exact
problem that was investigated is described in
Section~\ref{sec:ch1.problemStatement}.

%---------------------------------------------------------------------
%   Subject of investigation, quoted from the EEE4022S topic list
%   (RA Verrinder, 8 July 2026). Its references are listed below by hand,
%   not through BibTeX, so they stay separate from the main bibliography.
%---------------------------------------------------------------------
\section*{Subject of investigation}
\begin{quote}
Building on Zhao et al.\ (2022) and the surface characterisation work of Spirakis (2021); Spirakis et al.\ (2023);
Jones (2024), this project develops and validates a stereo-camera three-dimensional reconstruction pipeline on the
UCT BlueROV2 platform, as a first step towards reconstructing the underside of Antarctic sea ice. The work is
conducted in a controlled pool environment, where a submerged proxy target of known geometry stands in for the ice
underside. External fiducial markers provide the camera pose reference, so that the reconstruction is decoupled from
onboard navigation, and the recovered surface model is evaluated against the ground-truth geometry of the proxy. The
focus is on the geometric reconstruction itself: stereo calibration, handling of the underwater optical path,
multi-view fusion, and quantified reconstruction accuracy.
\end{quote}

\section*{References}
\begin{list}{}{\setlength{\leftmargin}{1.5em}\setlength{\itemindent}{-1.5em}\setlength{\itemsep}{0.3em}}
\item Blue Robotics. BlueROV2. [ONLINE], 2025. Accessed: 2025-07-11.
  url: \url{https://bluerobotics.com/store/rov/bluerov2/}
\item D. Jones. SHARC project: Ice floe size and shape characterisation from LiDAR scans. BSc(Eng) final year
  project, Department of Electrical Engineering), University of Cape Town, 2024.
  url: \url{https://tinyurl.com/djones2023}
\item J. Landy, A. Komarov, and D. Barber. Numerical and experimental evaluation of terrestrial LiDAR for
  parameterizing centimeter-scale sea ice surface roughness. \textit{IEEE Transactions on Geoscience and Remote
  Sensing}, 53:4887--4898, 2015. doi: \href{https://doi.org/10.1109/TGRS.2015.2412034}{10.1109/TGRS.2015.2412034}
\item J. Landy, A. Komarov, and D. Barber. Secondary-scale surface roughness parameterization using terrestrial
  LiDAR. In \textit{IGARSS 2015-2015 IEEE International Geoscience and Remote Sensing Symposium}, pages
  2068--2017. IEEE, 2023. doi: \href{https://doi.org/10.1109/igarss.2015.7326208}{10.1109/igarss.2015.7326208}
\item A. Spirakis. Investigation into surface roughness characterisation using LiDAR and cameras. BSc(Eng) final
  year project, Department of Electrical Engineering), University of Cape Town, 2021.
  url: \url{https://tinyurl.com/spirakis2021}
\item A. Spirakis, J. Hepworth, and R. Verrinder. 3D reconstruction of pancake sea ice using LiDAR and cameras. In
  \textit{IGARSS 2023-2023 IEEE International Geoscience and Remote Sensing Symposium}, pages 60--63. IEEE, 2023.
  doi: \href{https://doi.org/10.1109/igarss52108.2023.10282250}{10.1109/igarss52108.2023.10282250}
\item M. Vichi. An indicator of sea ice variability for the Antarctic marginal ice zone. \textit{The Cryosphere},
  16:4087--4106, 2022. doi: \href{https://doi.org/10.5194/tc-16-4087-2022}{10.5194/tc-16-4087-2022}
\item L. Zhao, M. Zhou, and B. Loose. Towards under-ice sensing using a portable ROV. In \textit{OCEANS
  2022-Hampton Roads}, pages 1--8. IEEE, 2022.
  doi: \href{https://doi.org/10.1109/OCEANS47191.2022.9977140}{10.1109/OCEANS47191.2022.9977140}
\end{list}
